\documentclass[aspectratio=169,11pt]{beamer}

% Vietnamese-compatible, pdfLaTeX-friendly setup
\usepackage[utf8]{inputenc}
\usepackage[T5]{fontenc}
\usepackage[vietnamese]{babel}

% Tables, figures, and diagrams
\usepackage{graphicx}
\usepackage{svg}
\usepackage{booktabs}
\usepackage{array}
\usepackage{tikz}
\usetikzlibrary{arrows.meta,positioning,shapes.geometric}

% Theme and colors
\definecolor{PrismNavy}{HTML}{123B5D}
\definecolor{PrismBlue}{HTML}{1677B8}
\definecolor{PrismCyan}{HTML}{28A7C7}
\definecolor{PrismGold}{HTML}{E6A23C}
\definecolor{PrismLight}{HTML}{EEF5F8}
\definecolor{PrismText}{HTML}{22313F}
\definecolor{PrismMuted}{HTML}{627482}

\usetheme{default}
\usefonttheme{professionalfonts}
\setbeamercolor{normal text}{fg=PrismText,bg=white}
\setbeamercolor{structure}{fg=PrismBlue}
\setbeamercolor{frametitle}{fg=white,bg=PrismNavy}
\setbeamercolor{title}{fg=PrismNavy}
\setbeamercolor{subtitle}{fg=PrismBlue}
\setbeamercolor{block title}{fg=white,bg=PrismBlue}
\setbeamercolor{block body}{fg=PrismText,bg=PrismLight}
\setbeamercolor{alerted text}{fg=PrismGold}
\setbeamercolor{footline}{fg=white,bg=PrismNavy}
\setbeamertemplate{navigation symbols}{}
\setbeamertemplate{itemize item}{\color{PrismCyan}$\blacktriangleright$}
\setbeamertemplate{itemize subitem}{\color{PrismBlue}$\bullet$}
\setbeamertemplate{enumerate item}{\color{PrismBlue}\insertenumlabel.}
\setbeamertemplate{frametitle}{%
  \nointerlineskip
  \begin{beamercolorbox}[wd=\paperwidth,ht=0.82cm,dp=0.22cm,leftskip=0.55cm]{frametitle}
    \usebeamerfont{frametitle}\insertframetitle
  \end{beamercolorbox}
}
\setbeamertemplate{footline}{%
  \begin{beamercolorbox}[wd=\paperwidth,ht=0.32cm,dp=0.13cm,leftskip=0.45cm,rightskip=0.45cm]{footline}
    \scriptsize Nhóm 8 -- Phân tích điểm thi THPT
    \hfill
    \insertframenumber/\inserttotalframenumber
  \end{beamercolorbox}
}
\setbeamersize{text margin left=0.55cm,text margin right=0.55cm}
\setlength{\tabcolsep}{5pt}
\renewcommand{\arraystretch}{1.12}
\svgsetup{inkscapelatex=false}

% Use the same image paths as the report. If an asset is temporarily absent,
% keep the slide compilable and show its intended path as a placeholder.
\newcommand{\MissingImage}[2][0.72\textheight]{%
  \begin{tikzpicture}
    \node[
      draw=PrismBlue,dashed,rounded corners=3pt,
      fill=PrismLight,minimum width=0.91\linewidth,
      minimum height=#1,align=center,text=PrismMuted
    ] {\Large Ảnh placeholder\\[0.25cm]\ttfamily\detokenize{#2}};
  \end{tikzpicture}%
}
\newcommand{\ReportImage}[2][0.72\textheight]{%
  \IfFileExists{#2}{\includegraphics[width=\linewidth,height=#1,keepaspectratio]{#2}}{\MissingImage[#1]{#2}}%
}
\newcommand{\ReportSVG}[2][0.72\textheight]{%
  \IfFileExists{#2}{\includesvg[width=\linewidth,height=#1,keepaspectratio]{#2}}{\MissingImage[#1]{#2}}%
}
\newcommand{\Metric}[3]{%
  \begin{beamercolorbox}[wd=0.29\linewidth,sep=0.18cm,center,rounded=true]{block body}
    {\color{PrismBlue}\Large\bfseries #1}\par
    {\scriptsize #2}\par
    {\tiny\color{PrismMuted}#3}
  \end{beamercolorbox}%
}

\title{Phân tích và trực quan hóa\\điểm thi tốt nghiệp THPT Việt Nam}
\subtitle{Dashboard tương tác và AI Assistant có kiểm soát}
\author{Nhóm 8}
\institute{Khoa Công nghệ Thông tin\\Trường Đại học Khoa học Tự nhiên -- ĐHQG-HCM}
\date{Tháng 7 năm 2026}

\begin{document}

\begin{frame}[plain]
  \begin{columns}[c,onlytextwidth]
    \begin{column}{0.25\textwidth}
      \centering
      \ReportImage[3.15cm]{figures/hcmus-logo.png}
    \end{column}
    \begin{column}{0.71\textwidth}
      {\small\color{PrismMuted}BÁO CÁO ĐỒ ÁN -- TRỰC QUAN HÓA DỮ LIỆU}\par
      \vspace{0.25cm}
      {\LARGE\bfseries\color{PrismNavy}\inserttitle}\par
      \vspace{0.18cm}
      {\large\color{PrismBlue}\insertsubtitle}\par
      \vspace{0.55cm}
      {\normalsize\insertauthor}\par
      {\small\color{PrismMuted}\insertinstitute}\par
      \vspace{0.25cm}
      {\small\insertdate}
    \end{column}
  \end{columns}
\end{frame}

\begin{frame}{Nhóm thực hiện}
  \small
  \begin{tabular}{@{}p{3.1cm}p{1.8cm}p{8.95cm}@{}}
    \toprule
    \textbf{Thành viên} & \textbf{MSSV} & \textbf{Phụ trách chính} \\
    \midrule
    Phạm Quốc Khánh & 23120283 & Thu thập và mô tả nguồn dữ liệu \\
    Phạm Thành Nam & 23120301 & Khảo sát, tiền xử lý và chuẩn hóa dữ liệu \\
    Trương Quang Phát & 23120318 & Dashboard frontend, tab Tổng quan và Xu hướng \\
    Châu Huỳnh Phúc & 23120329 & Phổ điểm, tổ hợp, địa phương và vùng miền \\
    Huỳnh Tấn Phước & 23120334 & AI Assistant, FastAPI, thực thi local và logs \\
    \bottomrule
  \end{tabular}
\end{frame}

\begin{frame}{Bài toán và mục tiêu}
  \begin{columns}[T,onlytextwidth]
    \begin{column}{0.48\textwidth}
      \begin{block}{Bài toán}
        \begin{itemize}
          \item Dữ liệu điểm thi có quy mô lớn và nhiều chiều phân tích.
          \item Bảng dữ liệu thô khó bộc lộ xu hướng, phân bố và khác biệt địa lý.
          \item CT2006 và CT2018 có cấu trúc môn học khác nhau.
        \end{itemize}
      \end{block}
    \end{column}
    \begin{column}{0.48\textwidth}
      \begin{block}{Mục tiêu}
        \begin{itemize}
          \item Xây dựng dashboard rõ ràng, có bộ lọc và điều hướng.
          \item Dẫn dắt từ tổng quan đến môn học, tổ hợp và địa phương.
          \item Hỗ trợ hỏi dữ liệu bằng ngôn ngữ tự nhiên qua AI Assistant.
        \end{itemize}
      \end{block}
    \end{column}
  \end{columns}
  \vspace{0.25cm}
  \centering
  \alert{Biến dữ liệu điểm thi thành một câu chuyện trực quan có thể khám phá và kiểm chứng.}
\end{frame}

\begin{frame}{Phạm vi và cấu trúc hệ thống}
  \centering
  \begin{tikzpicture}[
    node distance=0.42cm,
    box/.style={rounded corners=4pt,draw=PrismBlue,fill=PrismLight,
      minimum width=3.25cm,minimum height=1.25cm,align=center,font=\small},
    arr/.style={-{Latex[length=3mm]},very thick,draw=PrismCyan}
  ]
    \node[box] (raw) {Dữ liệu thô\\TXT/CSV/XLSX\\2022--2026};
    \node[box,right=of raw] (etl) {Pipeline\\làm sạch, chuẩn hóa\\và tạo đặc trưng};
    \node[box,right=of etl] (dash) {Dashboard\\5 góc nhìn\\tương tác};
    \node[box,right=of dash] (ai) {AI Assistant\\sinh code, phê duyệt\\và thực thi local};
    \draw[arr] (raw) -- (etl);
    \draw[arr] (etl) -- (dash);
    \draw[arr] (dash) -- (ai);
  \end{tikzpicture}
  \vspace{0.65cm}

  \Metric{2022--2026}{Giai đoạn dữ liệu}{5 năm thi}
  \hfill
  \Metric{5.436.558}{Số bản ghi sau xử lý}{mỗi hàng là một thí sinh}
  \hfill
  \Metric{44}{Số biến}{định danh, địa lý, điểm, ngoại ngữ và tổ hợp}
\end{frame}

\begin{frame}{Nguồn và quy mô dữ liệu}
  \begin{columns}[T,onlytextwidth]
    \begin{column}{0.55\textwidth}
      \small
      \begin{tabular}{@{}lr@{}}
        \toprule
        \textbf{Năm} & \textbf{Số hàng sau xử lý} \\
        \midrule
        2022 & 995.435 \\
        2023 & 1.017.584 \\
        2024 & 1.061.604 \\
        2025 & 1.153.072 \\
        2026 & 1.208.863 \\
        \midrule
        \textbf{Tổng} & \textbf{5.436.558} \\
        \bottomrule
      \end{tabular}
    \end{column}
    \begin{column}{0.41\textwidth}
      \begin{block}{Đầu vào}
        \small
        \begin{itemize}
          \item Các tệp điểm thi 2022--2026 được adapter hóa từ nguồn thô.
          \item Mỗi bản ghi tương ứng một thí sinh.
          \item Schema khác nhau theo năm và chương trình.
        \end{itemize}
      \end{block}
      \begin{block}{Đầu ra}
        \small
        \texttt{data/processed/final\_data.csv}
      \end{block}
    \end{column}
  \end{columns}
\end{frame}

\begin{frame}{Pipeline tiền xử lý}
  \small
  \begin{columns}[T,onlytextwidth]
    \begin{column}{0.32\textwidth}
      \begin{block}{1. Chuẩn hóa}
        \begin{itemize}
          \item Đọc CSV/XLSX
          \item Đồng nhất tên cột
          \item Chuẩn hóa số báo danh
          \item Chuyển điểm sang kiểu số
        \end{itemize}
      \end{block}
    \end{column}
    \begin{column}{0.32\textwidth}
      \begin{block}{2. Làm sạch}
        \begin{itemize}
          \item Giữ điểm thiếu là \texttt{NaN}
          \item Loại điểm ngoài $[0,10]$
          \item Loại hàng không có điểm hợp lệ
          \item Ánh xạ tỉnh và vùng miền
        \end{itemize}
      \end{block}
    \end{column}
    \begin{column}{0.32\textwidth}
      \begin{block}{3. Tạo đặc trưng}
        \begin{itemize}
          \item Số môn và điểm trung bình
          \item Điểm ngoại ngữ theo mã N1--N7
          \item Ban/chương trình học
          \item 12 tổ hợp xét tuyển
        \end{itemize}
      \end{block}
    \end{column}
  \end{columns}
  \vspace{0.3cm}
  \centering
  \alert{Mục tiêu: một schema thống nhất, đủ tin cậy để mọi biểu đồ dùng chung.}
\end{frame}

\begin{frame}{Schema dữ liệu sau xử lý}
  \small
  \begin{tabular}{@{}p{3.1cm}p{11.5cm}@{}}
    \toprule
    \textbf{Nhóm biến} & \textbf{Các trường tiêu biểu} \\
    \midrule
    Định danh & \texttt{nam}, \texttt{chuong\_trinh}, \texttt{sbd} \\
    Địa lý & \texttt{ma\_tinh}, \texttt{ten\_tinh}, \texttt{vung\_mien}, \texttt{vung\_3} \\
    Điểm môn học & Toán, Ngữ văn, Ngoại ngữ, Vật lý, Hóa học, Sinh học, Lịch sử, Địa lý, GDCD, Tin học, Công nghệ, GDKTPL \\
    Biến dẫn xuất & \texttt{so\_mon}, \texttt{diem\_tb}, \texttt{ban}, \texttt{ten\_ngoai\_ngu}, các cột điểm ngoại ngữ riêng \\
    Điểm tổ hợp & \texttt{diem\_khoi\_a00} đến \texttt{diem\_khoi\_d15} cho 12 tổ hợp \\
    \bottomrule
  \end{tabular}
  \vspace{0.45cm}
  \begin{block}{Quy tắc quan trọng}
    Điểm tổ hợp chỉ được tính khi thí sinh có đủ môn thành phần; các biểu đồ dùng tên cột canonical, ví dụ \texttt{ngu\_van}, \texttt{vat\_li}, \texttt{diem\_khoi\_a00}.
  \end{block}
\end{frame}

\begin{frame}{Dashboard: năm góc nhìn liên kết}
  \begin{columns}[T,onlytextwidth]
    \begin{column}{0.48\textwidth}
      \begin{enumerate}
        \item \textbf{Tổng quan}: KPI và bức tranh toàn quốc
        \item \textbf{Xu hướng \& Môn học}: biến động theo thời gian
        \item \textbf{Phổ điểm \& Tổ hợp}: hình dạng phân bố
        \item \textbf{Địa phương \& Vùng miền}: xếp hạng và so sánh
        \item \textbf{Trợ lý AI}: hỏi, sinh code và kiểm duyệt
      \end{enumerate}
    \end{column}
    \begin{column}{0.48\textwidth}
      \begin{block}{Bộ lọc dùng chung}
        \begin{itemize}
          \item Năm và chương trình
          \item Môn học, chỉ số và tổ hợp
          \item Vùng miền, tỉnh/thành phố
          \item Số lượng phần tử Top-$N$
        \end{itemize}
      \end{block}
    \end{column}
  \end{columns}
  \vfill
  \centering
  \alert{Luồng khám phá: Toàn quốc $\rightarrow$ xu hướng $\rightarrow$ phân bố $\rightarrow$ địa phương.}
\end{frame}

\begin{frame}{Tab Tổng quan}
  \centering
  \ReportImage[6.65cm]{figures/overview_tab.png}
  \vspace{-0.12cm}

  {\scriptsize\color{PrismMuted}KPI quy mô, điểm trung bình và xu hướng toàn quốc theo trạng thái lọc.}
\end{frame}

\begin{frame}{Tab Xu hướng \& Môn học}
  \centering
  \ReportImage[6.65cm]{figures/subject_trend_tab.png}
  \vspace{-0.12cm}

  {\scriptsize\color{PrismMuted}Biểu đồ đường và heatmap hỗ trợ so sánh môn học theo năm/chương trình.}
\end{frame}

\begin{frame}{Tab Phổ điểm \& Tổ hợp}
  \centering
  \ReportImage[6.65cm]{figures/distribution_tab.png}
  \vspace{-0.12cm}

  {\scriptsize\color{PrismMuted}Histogram, thống kê tóm tắt và so sánh các tổ hợp xét tuyển.}
\end{frame}

\begin{frame}{Tab Địa phương \& Vùng miền}
  \centering
  \ReportImage[6.65cm]{figures/region_tab.png}
  \vspace{-0.12cm}

  {\scriptsize\color{PrismMuted}Top/Bottom tỉnh, điểm vùng và heatmap vùng $\times$ môn học.}
\end{frame}

\begin{frame}{Câu hỏi 1: Điểm trung bình tăng do đâu?}
  \begin{columns}[c,onlytextwidth]
    \begin{column}{0.62\textwidth}
      \centering
      \ReportSVG[5.55cm]{figures/q1_national_average_ct2006.svg}
    \end{column}
    \begin{column}{0.35\textwidth}
      \begin{block}{Phạm vi}
        \small
        \begin{itemize}
          \item CT2006
          \item Giai đoạn 2022--2024
          \item So sánh toàn quốc với 9 môn
        \end{itemize}
      \end{block}
      \centering
      {\color{PrismBlue}\LARGE\bfseries 6,40 $\rightarrow$ 6,75}\par
      \vspace{0.12cm}
      {\small Tăng \textbf{0,35 điểm}}
    \end{column}
  \end{columns}
\end{frame}

\begin{frame}{Câu hỏi 1: Xu hướng theo từng môn}
  \begin{columns}[c,onlytextwidth]
    \begin{column}{0.67\textwidth}
      \centering
      \ReportSVG[5.65cm]{figures/q1_subject_trends_ct2006.svg}
    \end{column}
    \begin{column}{0.30\textwidth}
      \small
      \begin{itemize}
        \item \textbf{Sinh học}: $+1,26$
        \item \textbf{Ngữ văn}: $+0,72$
        \item \textbf{Địa lý}: $+0,51$
        \item Tiếng Anh: $+0,36$
        \item Toán, Hóa, Lý: giảm nhẹ
      \end{itemize}
      \vspace{0.2cm}
      \begin{beamercolorbox}[sep=0.18cm,rounded=true]{block body}
        \centering \textbf{6/9 môn tăng}\par
        \scriptsize nhưng mức tăng không đồng đều
      \end{beamercolorbox}
    \end{column}
  \end{columns}
\end{frame}

\begin{frame}{Kết luận phân tích 1}
  \begin{block}{Phát hiện chính}
    Điểm trung bình toàn quốc tăng \textbf{không đồng nghĩa toàn bộ môn thi cùng cải thiện}. Xu hướng tăng chủ yếu tập trung ở \textbf{Sinh học, Ngữ văn và Địa lý}.
  \end{block}
  \vspace{0.3cm}
  \begin{columns}[T,onlytextwidth]
    \begin{column}{0.48\textwidth}
      \centering
      {\color{PrismBlue}\Huge\bfseries 0,35}\par
      Mức tăng toàn quốc
    \end{column}
    \begin{column}{0.48\textwidth}
      \centering
      {\color{PrismGold}\Huge\bfseries 1,31}\par
      Khoảng cách giữa môn tăng nhiều nhất\par và môn giảm nhiều nhất
    \end{column}
  \end{columns}
  \vfill
  \centering
  \alert{Cần luôn đặt chỉ số tổng hợp cạnh chuỗi điểm của từng môn.}
\end{frame}

\begin{frame}{Câu hỏi 2: Mặt bằng vùng hay cực trị địa phương?}
  \begin{columns}[c,onlytextwidth]
    \begin{column}{0.68\textwidth}
      \centering
      \ReportSVG[5.7cm]{figures/q2_region_snapshot_2026.svg}
    \end{column}
    \begin{column}{0.29\textwidth}
      \begin{block}{Trạng thái lọc}
        \small
        \begin{itemize}
          \item Năm 2026
          \item Môn Tổng hợp
          \item Phạm vi Toàn quốc
          \item Top 10 tỉnh/thành
        \end{itemize}
      \end{block}
      \small
      Đồng bằng sông Hồng dẫn đầu với khoảng \textbf{6,15 điểm}.
    \end{column}
  \end{columns}
\end{frame}

\begin{frame}{Câu hỏi 2: Kiểm tra độ phân hóa nội vùng}
  \begin{columns}[c,onlytextwidth]
    \begin{column}{0.68\textwidth}
      \centering
      \ReportSVG[5.7cm]{figures/q2_bottom_region_snapshot_2026.svg}
    \end{column}
    \begin{column}{0.29\textwidth}
      \small
      \begin{itemize}
        \item ĐBSH: khoảng cách nội vùng \textbf{0,48}
        \item Trung du \& MNPB: \textbf{0,93}
        \item Bắc Trung Bộ \& DHMT: \textbf{0,61}
      \end{itemize}
      \vspace{0.2cm}
      \begin{beamercolorbox}[sep=0.18cm,rounded=true]{block body}
        Chênh lệch nội vùng lớn nhất gần \textbf{gấp đôi} chênh lệch giữa vùng cao nhất và thấp nhất.
      \end{beamercolorbox}
    \end{column}
  \end{columns}
\end{frame}

\begin{frame}{Kết luận phân tích 2}
  \begin{block}{Phát hiện chính}
    Chênh lệch vùng miền \textbf{không được tạo ra bởi mọi địa phương theo cùng một mức độ}. Điểm trung bình vùng có thể che khuất sự phân hóa lớn giữa các tỉnh trong cùng vùng.
  \end{block}
  \vspace{0.35cm}
  \begin{columns}[T,onlytextwidth]
    \begin{column}{0.31\textwidth}
      \centering
      {\color{PrismBlue}\Large\bfseries Mặt bằng vùng}\par
      \small Xác định vị trí tương đối
    \end{column}
    \begin{column}{0.31\textwidth}
      \centering
      {\color{PrismCyan}\Large\bfseries Top/Bottom tỉnh}\par
      \small Nhận diện các cực trị
    \end{column}
    \begin{column}{0.31\textwidth}
      \centering
      {\color{PrismGold}\Large\bfseries Khoảng cách nội vùng}\par
      \small Đánh giá mức độ đồng đều
    \end{column}
  \end{columns}
  \vfill
  \centering
  \alert{Ba lớp thông tin cần được đọc đồng thời khi so sánh địa lý.}
\end{frame}

\begin{frame}{Kiến trúc AI Assistant}
  \begin{columns}[c,onlytextwidth]
    \begin{column}{0.69\textwidth}
      \centering
      \ReportImage[5.8cm]{figures/system_architecture.png}
    \end{column}
    \begin{column}{0.28\textwidth}
      \small
      \begin{itemize}
        \item Frontend nhận câu hỏi
        \item AI API sinh code và giải thích
        \item Execution API validate và chạy local
        \item Attachments API phân tích file đính kèm
        \item Logs/Report API lưu và tổng hợp lịch sử
      \end{itemize}
      \vspace{0.2cm}
      \begin{beamercolorbox}[sep=0.18cm,rounded=true]{block body}
        \centering \textbf{Sinh code và thực thi là hai bước tách biệt.}
      \end{beamercolorbox}
    \end{column}
  \end{columns}
\end{frame}

\begin{frame}{AI Assistant: từ câu hỏi đến code đề xuất}
  \begin{columns}[c,onlytextwidth]
    \begin{column}{0.68\textwidth}
      \centering
      \ReportImage[5.8cm]{figures/ai_assistant_tab.png}
    \end{column}
    \begin{column}{0.29\textwidth}
      \small
      \begin{enumerate}
        \item Người dùng nhập câu hỏi hoặc đính kèm file.
        \item AI sinh code, giải thích và cảnh báo nếu có.
        \item Code ở trạng thái chờ kiểm tra.
      \end{enumerate}
      \vspace{0.25cm}
      \alert{Không có thực thi ngầm.}
    \end{column}
  \end{columns}
\end{frame}

\begin{frame}{AI Assistant: phê duyệt và thực thi local}
  \begin{columns}[T,onlytextwidth]
    \begin{column}{0.49\textwidth}
      \centering
      \ReportImage[4.65cm]{figures/ai_code_review.png}
      \par\scriptsize\color{PrismMuted}Người dùng xem, chỉnh sửa, từ chối hoặc phê duyệt code.
    \end{column}
    \begin{column}{0.49\textwidth}
      \centering
      \ReportImage[4.65cm]{figures/ai_execution_result.png}
      \par\scriptsize\color{PrismMuted}Chỉ phiên bản đã phê duyệt mới được chạy cục bộ.
    \end{column}
  \end{columns}
  \vfill
  \centering
  \alert{JSON logs lưu yêu cầu, code ban đầu, code cuối, kết quả, lỗi và ảnh biểu đồ để truy vết.}
\end{frame}

\begin{frame}{Đánh giá theo yêu cầu môn học}
  \small
  \begin{tabular}{@{}p{11.8cm}>{\centering\arraybackslash}p{2cm}@{}}
    \toprule
    \textbf{Tiêu chí} & \textbf{Trạng thái} \\
    \midrule
    Dữ liệu thực tế liên quan đến Việt Nam & \color{PrismBlue}\textbf{Đạt} \\
    Dữ liệu trên 2.000 hàng và có ít nhất 7 biến & \color{PrismBlue}\textbf{Đạt} \\
    Ứng dụng trực quan hóa dựa trên dashboard & \color{PrismBlue}\textbf{Đạt} \\
    Có bộ lọc tương tác và điều hướng & \color{PrismBlue}\textbf{Đạt} \\
    Có câu chuyện phân tích trực quan & \color{PrismBlue}\textbf{Đạt} \\
    Có tích hợp AI với cơ chế kiểm soát & \color{PrismBlue}\textbf{Đạt} \\
    \bottomrule
  \end{tabular}
  \vspace{0.4cm}
  \begin{beamercolorbox}[sep=0.2cm,rounded=true]{block body}
    Trạng thái cần được đối chiếu lần cuối với phiên bản dashboard và minh chứng nộp chính thức.
  \end{beamercolorbox}
\end{frame}

\begin{frame}{Hạn chế và hướng phát triển}
  \begin{columns}[T,onlytextwidth]
    \begin{column}{0.48\textwidth}
      \begin{block}{Hạn chế}
        \small
        \begin{itemize}
          \item Khác biệt CT2006 và CT2018.
          \item Ánh xạ địa phương sau sáp nhập.
        \item Tệp dữ liệu cuối khoảng 785 MB và không commit lên git.
          \item Kết quả AI vẫn có thể sai logic.
        \end{itemize}
      \end{block}
    \end{column}
    \begin{column}{0.48\textwidth}
      \begin{block}{Hướng phát triển}
        \small
        \begin{itemize}
          \item Hoàn thiện API dữ liệu tổng hợp theo bộ lọc.
          \item Chuyển logs từ JSON sang SQLite/Prisma khi cần.
          \item Trực quan hóa bản đồ địa lý.
          \item Mẫu câu hỏi AI và kiểm thử tự động.
          \item Xuất dữ liệu và báo cáo.
        \end{itemize}
      \end{block}
    \end{column}
  \end{columns}
\end{frame}

\begin{frame}{Kết luận}
  \begin{columns}[c,onlytextwidth]
    \begin{column}{0.25\textwidth}
      \centering
      \ReportImage[2.65cm]{figures/hcmus-logo.png}
    \end{column}
    \begin{column}{0.70\textwidth}
      \begin{itemize}
        \item Hợp nhất \textbf{5,44 triệu} bản ghi điểm thi giai đoạn 2022--2026.
        \item Xây dựng dashboard gồm tổng quan, xu hướng, phân bố và địa lý.
        \item Hai phân tích cho thấy chỉ số tổng hợp có thể che khuất khác biệt ở cấp môn và địa phương.
        \item AI Assistant mở rộng khả năng khám phá dữ liệu nhưng giữ quyền kiểm soát ở người dùng.
      \end{itemize}
      \vspace{0.45cm}
      {\Large\bfseries\color{PrismNavy}Xin cảm ơn!}\par
      \vspace{0.12cm}
      {\color{PrismBlue}Q\&A}
    \end{column}
  \end{columns}
\end{frame}

\end{document}
